\documentclass[runningheads]{llncs}
\usepackage{longtable}
\usepackage{nth}
\usepackage[english=american,babel=false]{csquotes}

\begin{document}
\frontmatter
\setcounter{page}{5}

\chapter*{Preface to the Proceedings of HS3 2026}
\markboth{Preface}{Preface}

This volume presents the revised papers of the \nth{2} Workshop on
Hardware-Supported Software Security (HS3 2026), co-located with the
\nth{31} European Symposium on Research in Computer Security (ESORICS
2026), held in Rome, Italy, on September 17, 2026.

The HS3 workshop seeks to share experience, tools and methodology on
hardware-assisted software security. Combining software and hardware
aspects to mitigate new software attacks is becoming increasingly
important. For example, hardware vulnerabilities such as Spectre or
Meltdown can be exploited entirely through software-level attacks, and these
attacks can be executed remotely, without requiring physical access to the
targeted hardware platform. Effective mitigations to these attacks require
hardware/software co-design, in particular when performance impacts are to
be minimized. In addition, hardware features can also be used to better
detect and respond to traditional software attacks, such as memory
corruption. The goal of the HS3 workshop is to provide a forum for
researchers and practitioners from academia, industry and government who
study the security of software/hardware interfaces in terms of attacks and
defenses to develop hardware-supported software security.

For HS3 2026 we were looking for submissions that propose new processor
architectures or processor extensions that offer better resilience against
software attacks. These architectures typically provide hardware-based
security mechanisms to protect the software stack. One of the key
challenges in the field is to formally specify and verify the security
guarantees offered by such architectures and to better assess the security
guarantees provided by existing hardware architectures against software
attacks, especially attacks against micro-architecture. This can be
achieved by identifying new vulnerabilities using reverse engineering,
fuzzing or other attack approaches. HS3 2026 also had a special theme on
\enquote{Hardware-Supported Software Security for AI Systems} and
we were specifically inviting submissions that present hardware-assisted
approaches in this domain.

Our call for papers attracted significant interest and resulted in 20
submissions. Each manuscript underwent a rigorous peer-review process, with
three single-blinded reviewers per paper, and guaranteeing that chairs and
committee members would not review their own or otherwise conflicting
submissions. Ultimately, we accepted eight regular papers with an
acceptance rate of 40\,\%. We further accepted three submissions as
work-in-progress presentations where the program committee felt that the
material was interesting fuel for discussions while the overall work was not
yet sufficiently complete to warrant publication.  The workshop itself then
featured 11 presentations and time for discussion and reflection. With
about 30 participants, the workshop was well attended and promoted lively
discussions around recent trends in hardware-supported security and on
application-related topics.

Our thanks go to the program committee and external reviewers for their
tireless commitment during the advertising and reviewing phases.  Most
importantly, we thank all our speakers and participants for the great
presentations and intriguing discussions during the workshop. Finally, we
would like to thank the ESORICS organizers and specifically the workshop
chairs for supporting us throughout the process of preparing and running
HS3.

\medskip

\begin{flushright}
\noindent October 2026\hfill
{Thomas Rokicki} \\
{Jan Tobias M\"uhlberg} \\
{Yuko Hara} \\
\end{flushright}


\chapter*{Organization of HS3 2026}
\markboth{Organization}{Organization}
\setlength{\tabcolsep}{0pt}     % remove spacing between columns in the following long tables
\setlength{\LTpre}{0pt}         % remove extra vertical space before the long tables

\section*{Program Committee Chairs}
\begin{longtable}{p{0.3\textwidth}p{0.7\textwidth}}
Thomas Rokicki       & CentraleSup\'elec/Inria, France \\
Jan Tobias M\"uhlberg & Universit\'e Libre de Bruxelles, Belgium \\
Yuko Hara            & CNRS, France \\
\end{longtable}

\section*{Program Committee}
\begin{longtable}{p{0.3\textwidth}p{0.7\textwidth}}
Iness Ben Guirat     & Universit\'e Libre de Bruxelles, Belgium \\
Kevin Cheang         & University of California, Berkeley, USA  \\
Pascal Cotret        & ENSTA Bretagne, France \\
Chris Dalton         & HP Labs, United Kingdom \\
Lesly-Ann Daniel     & EURECOM, France \\
Merve G\"ulmez       & Ericsson Research, Sweden \& KU Leuven, Belgium \\
Karine Heydemann     & Thales, France \\
Guillaume Hiet       & IRISA, CentraleSup\'elec/Inria, France \\
Vianney Lap\^otre    & Universit\'e Bretagne Sud, France \\
Cl\'ementine Maurice & CNRS, France \\
Maria M\'endez Real  & Universit\'e Bretagne Sud, France \\
Cristofaro Mune      & Raelize B.V., Malta \\
Antonio Mu\~noz      & University of M\'alaga, Spain \\
Kaveh Razavi         & ETH Zurich, Switzerland \\
Benedict Schl\"uter  & ETH Zurich, Switzerland \\
Volker Stolz         & HVL, Norway \\
Marcus V\"olp        & University of  Luxembourg, Luxembourg \\
Pierre Wilke         & CentraleSup\'elec/Inria, France \\
\end{longtable}

\section*{Additional Reviewers}
\begin{multicols}{2}
\noindent Navid Ladner \\
\end{multicols}

\tableofcontents

\mainmatter

%Create dummy chapters for the sample ToC

\setcounter{page}{1}
\addtocontents{toc}{\protect\section*{Side-Channel Attacks \& Discovery}}
\author{Eyal Hadad and Mordechai Guri}
\title{Shape and Substance: Dual-Layer Side-Channel Attacks on Local Vision-Language Models}
\maketitle
\clearpage

\setcounter{page}{2}
\author{Abla Smahi and Jan Tobias M\"uhlberg}
\title{SoK: Software-Controlled Side Channels in LLM Inference}
\maketitle
\clearpage

\setcounter{page}{3}
\author{Nathana\"el Simon, Maria Mushtaq and Ludovic Apvrille}
\title{RLeak: Hybrid Reinforcement Learning Guided Fuzzing Framework for RISC-V Vulnerability Timing Side-Channel Discovery}
\maketitle
\clearpage

\setcounter{page}{4}
\addtocontents{toc}{\protect\section*{Measurement, Detection \& Evaluation Methodologies}}
\author{Adam Henault, Camille Moni\`ere, Philippe Tanguy and Vianney Lap\^otre}
\title{Statistical Analysis of HPC Filtering: Towards Reliable Use in Security}
\maketitle
\clearpage

\setcounter{page}{5}
\author{Cl\'ement M\'edart, Pierre Graux, Cl\'ementine Maurice and Gilles Grimaud}
\title{PerfCT: Detecting Constant-Time Violations Using Hardware Performance Counters}
\maketitle
\clearpage

\setcounter{page}{6}
\author{Damien Jauvart and Aur\'elien Vasselle}
\title{Evaluating the Security of Open-Source Cryptographic Libraries Against Physical Attacks: An End-to-End Study on Botan}
\maketitle
\clearpage

\setcounter{page}{7}
\addtocontents{toc}{\protect\section*{Trust \& Cryptographic Hardware}}
\author{Micha{\l} Wro\'nski, {\L}ukasz Dzierzkowski, Mateusz Le\'sniak and Ewa Syta}
\title{A Universal Weierstrass Engine for Hardware-Supported Quantum Cryptanalysis of RSA, DH, and ECC}
\maketitle
\clearpage

\setcounter{page}{8}
\author{Julian Funk, Matti Schulze, Patrick Sieber, Manuel V\"ogele, Jonas
R\"ockl and Tilo M\"uller}
\title{Solving Attestation Fragility: A Software-Based Trust Anchor for Agile Confidential VMs}
\maketitle

\clearpage
\setcounter{page}{9}
\phantom{Author Index}
\addcontentsline{toc}{title}{\protect\textbf{Author Index}}
\clearpage

\end{document}

